\section{Marked-root geometry}

For $(a,b,c)\in\A^3$, consider the binary cubic
\[
Q_{a,b,c}(U,V)=cU^3-2U^2V+bUV^2-2aV^3.
\]
Let $I\subset\A^3\times\Pone$ be its root incidence and let
$I^{\mathrm{simp}}$ denote the locus where the marked root is simple.

\begin{theorem}\label{thm:marked-root}
The morphism
\[
\A^3\longrightarrow I^{\mathrm{simp}},\qquad
(x,y,z)\longmapsto(F(x,y,z),[1+xy:x])
\]
is an isomorphism. Under this isomorphism, $F$ is the finite projection which
forgets the marked simple root.
\end{theorem}

On $V\ne0$, put $T=U/V$ and
$P(T)=cT^3-2T^2+bT-2a$.  A simple root $t$ reconstructs by
\[
v=P'(t)/2,qquad x=v^{-1},qquad y=t-v,qquad
z=5v^2-3tv-cv^3.
\]
The complementary $U\ne0$ chart is regular at the projective root at infinity
and supplies the $x=0$ source chart.  The two formulas agree on their overlap.
This second chart is essential: an affine-root computation alone would miss
part of the source.

Intrinsically, this is the restriction of
\[
\Pone\times\operatorname{Sym}^2(\Pone)
\longrightarrow\operatorname{Sym}^3(\Pone)
\]
after removing the repeated-marked-root divisor and the inverse image of a
tangent, nonosculating hyperplane.
